\section{Discussion}
\label{sec:discussion}

\subsection{Pre-execution commitment is not task typing}
\label{sec:discussion:not_typing}

The closest neighbouring frame, task-aware delegation~\cite{taskaware2026}, types a request and routes to whichever agent historically wins on the task type. The empirical evidence we report in Section~\ref{sec:results:exp1} suggests that under task typing the same delegation can route to the same agent and yet have its operational responsibility structure resolved very differently across model families. Median cross-family Jaccard distance on the active set is comparable in magnitude to median cosine distance on the weight vector, indicating that families disagree not only about how strongly each dimension matters but also about which dimensions are in scope at all. Routing alone cannot fix this: the agent that arrives at the task carries its own projection. Pre-execution commitment is therefore a layer below routing, and visible only when projections from multiple agents are elicited and compared on the same delegation.

\subsection{Clarification has cost; the trigger should respect it}
\label{sec:discussion:clarification}

A naive ``always clarify on ambiguity'' policy is dominated by direct execution on simple delegations and dominated by clarification on truly ambiguous ones, and the long-horizon underspecification literature documents the cost~\cite{lhaw2024}. Our composite trigger (Equation~(\ref{eq:trigger})) folds three signals --- bidder-level \texttt{clarify}, projection divergence $D_\pi$, and inter-bidder coverage variance --- so that clarification activates when the joint signal is high relative to its cost. The pilot evidence under the v1.3 Anthropic-excluded cross panel that all $5$ of $5$ ambiguous-class examples elicit at least one family to clarify (Section~\ref{sec:results:exp1:clarification}) shows that the trigger is no longer conservative under v1.3, but it also fires on $58$--$86\%$ of all pilot examples regardless of class --- the opposite calibration regime from the v1.0 panel, where the legacy \texttt{gpt-4o-mini} never clarified and only $1/5$ ambiguous-class examples elicited any family. The engineering question for the main run is therefore not whether to add a trigger but how to calibrate $\tau_D$ and $\tau_V$ so that the trigger fires on the cases where it would help and stays silent on cases where it would not. This is a calibration problem under both panel regimes, not a presence-of-clarification problem.

\subsection{Settlement evidence reliability is a separate axis}
\label{sec:discussion:evidence}

Our settlement step (Section~\ref{sec:protocol:settlement}) treats evidence sources as having different reliabilities, with deterministic and harness evidence ranked above human, LLM-judge, and retrieval-based evidence (Equation~(\ref{eq:kappa})). The initial pilot used a three-judge LLM panel for the annotation reliability track (\S\ref{sec:results:dataset:annotation}); the redesigned measurement layer of the settlement loss (\S\ref{sec:results:exp2}) uses a 12-judge tier-stratified Anthropic-excluded panel, since deterministic verification is rarely available for paper-section quality. The R1.7 annotation reliability failure ($\alpha = 0.219$; Section~\ref{sec:results:dataset:alpha}) is the cleanest illustration: citation accuracy is the dimension where blind LLM evidence is weakest, and where retrieval-based evidence (literature search) and human evidence would help most. The main run should treat evidence-source reliability as a deployment knob rather than a fixed setting, and prioritize deterministic and harness evidence wherever it can be constructed (e.g., automatic citation-coverage checks against a fixed corpus). Section~\ref{sec:discussion:specifiability} discusses why \emph{human-evidence reliability itself} also has a binding ceiling that is not solved by adding raters, based on the boundary-condition finding in Section~\ref{sec:results:limitations:specifiability}.

\subsection{Why the closure of the taxonomy matters}
\label{sec:discussion:closure}

We observed empirically that LLM annotators differ widely in their open-vocabulary ``responsibility labels'' for the same artifact, and that this disagreement collapses when projection is constrained to a closed weight vector over a fixed dimension index. Closed-taxonomy projection is what makes Experiment 1's metric well-defined; without closure, $1 - \cos(r_A, r_B)$ would be undefined when $r_A$ and $r_B$ live in different label spaces. The closure also enables the active-set rule, the under- and over-projection metrics, and the per-dimension settlement loss --- all of which require a fixed dimension index. The cost of closure is that any case the taxonomy does not cover is either projected onto the nearest dimension (information loss) or excluded (coverage loss). $J_{v1.1}$'s $|J| = 12$ is a deliberate first-cycle commitment; expansion to $|J| = 38$ is reserved for future cycles.

\subsection{Why this paper studies R1 only}
\label{sec:discussion:r1_only}

The long-term plan covers five delegation categories (R1--R5) and 38 dimensions. The first cycle is restricted to R1 (paper-research delegation) with seven category dimensions because (i) paper writing supplies a controlled artifact distribution where load-bearing flaws can be engineered cleanly, (ii) the responsibility dimensions are stable across our research community, and (iii) the cost of pilot-scale annotation by LLM judges is low. The category-conditional results from R1 do not transfer automatically to R2 (code review and bug-fix), which would require a different evidence model (deterministic test execution, type checks, runtime errors) and a different responsibility ontology (correctness, security, test coverage, style). We treat R2--R5 as future work and disclose the scope throughout.

\subsection{What the within-model baseline does and does not rule out}
\label{sec:discussion:within_model}

Within-model variance at $T = 0.5$ is a stochasticity upper bound conditional on a fixed model and a fixed prompt. It does not bound (i) cross-temperature variance for the same model, (ii) cross-prompt variance under semantically equivalent prompts, or (iii) cross-version variance as the same provider updates the same nominal model name. The intentional asymmetry of cross-family at $T = 0$ versus within-model at $T = 0.5$ is conservative against (i), but the paper's empirical claims do not address (ii) or (iii). Both are concerns that the main run should pre-register: prompt sensitivity through paraphrase ablation, and version stability through repeated runs across model snapshots.

\subsection{Pilot annotation versus main-run annotation}
\label{sec:discussion:annotation}

Pass-2 LLM annotation reaches $\alpha \geq 0.4$ on four of seven R1 dimensions, leaves R1.1 and R1.3 borderline, and fails on R1.7. The failure on R1.7 is interpretable as a structural limit of blind LLM annotation: citation accuracy requires recall of the relevant literature's canonical works, which mid-tier LLMs do not consistently have. The borderline cases on R1.1 and R1.3 may resolve under human annotation; the failure on R1.7 will require either human annotators or retrieval-augmented annotation pipelines. Until that step is taken, R1.7 results in this paper carry an annotation-uncertainty caveat (Section~\ref{sec:results:limitations}); R1.1 and R1.3 carry a softer caveat.

The pilot's planned human-anchor extension on R1.1 / R1.4 / R1.7 was attempted, partially executed, and superseded by a methodological boundary-condition finding (Section~\ref{sec:results:limitations:specifiability}). The closed substitute --- a Prolific R1.1 salvage with three validated raters at a 33\% pass rate, plus a closed R1.7 author + peer co-annotation pass on six citation-rich packages --- is reported in place of breadth validation. Section~\ref{sec:discussion:specifiability} discusses the implication: human anchor at scale is bound by anchor specifiability and domain-expertise gating, not rater count.

\subsection{Anchor specifiability is a binding constraint on human anchor scalability}
\label{sec:discussion:specifiability}

The pilot's planned 270-rating Prolific extension treated rater throughput as the binding factor and budgeted accordingly. Two empirical observations undercut that assumption.

First, even with a sharpened citation-event-count anchor printed at the top of every form page and a required \texttt{event\_count} field, two independent peers in a pre-protocol pass produced score--event mappings decoupled from the anchor table: their scores plateaued at 4.0--4.5 across submitted event counts 2--6. On a definitional zero-event stress-test package whose agent output explicitly disclaims citation engagement, both peers counted non-citation actions as events and scored s=4.0--4.5. The form-embedded anchor, sharpened though it was, did not transfer the anchor's intent to peers reading it for the first time mid-session.

Second, two peers in a post-protocol pass --- having received a separately-distributed rater protocol document and re-evaluated the same six packages --- produced monotonic score--event mappings, both passed the zero-event stress test, and (in one case) reached the anchor table's $s = 5$ row that no pre-protocol rater reached. The pre/post-protocol contrast at $n = 2$ per group is not inferential, but the pattern is consistent: a separately-read protocol document appears to help in this small pass, and the form-embedded layer alone does not.

The binding constraint on human anchor in this setting is therefore \emph{specification specifiability + domain-expertise gating}, not rater throughput. Two layered specifiability constraints are partially closable by a separately-read protocol document --- counting convention (cluster vs.\ item) and score--event mapping --- but a third constraint (per-event correctness verification) requires source-of-truth domain knowledge and is decoupled from form-design levers. We document the operational response in a rater protocol (two-axis scoring: Axis 1 engagement count, applicable by any rater; Axis 2 correctness, expert-gated and sparse) and report it as a methodological finding rather than a remediated weakness. The implication for any project that proposes to use crowdworker human anchor as the validity backstop for an LLM-only annotator panel is that the binding cost is not per-rater but per-rater-protocol-design-iteration, and that the rater training step must be treated as a distinct prerequisite to the rating session, not embedded in the form.

Residual gaps after protocol-reading remain --- the post-protocol raters disagreed on a missing-comparator-cluster package, validating the protocol's own list of ``borderline cases v1 does not yet resolve.'' This is consistent with the broader claim of the section: specifiability is a binding ceiling that closes incrementally with each protocol revision, not a throughput problem.

\subsection{Deployment implications}
\label{sec:discussion:deployment}

Two deployment lessons follow from the pilot. First, an earlier qualitative reading recommended turning projection-driven execution on by default for prescriptive single-dim tasks; the headline pilot result reverses the directional advantage (Section~\ref{sec:results:exp2:aggregate}) and shows the cost concentrates on dims requiring deep specialty work --- R1.4 novelty mapping and R1.7 citation audit. Default-on is therefore not justified by the pilot; deployment should treat projection guidance as a dim-conditional move. Second, R1.7 (where Stage 1 directly applies) and R1.4 (where the cross-link is inferential, since R1.4 was not part of the Stage 1 closure) are the dims for which projection guidance hurts most. The protocol-design moves that may help these dims --- retrieval-augmentation for R1.7, a bounded prior-work corpus and structured comparator template for R1.4 --- are dim-specific and not addressed by generic active-set or format hygiene. The active-set tightening that an earlier hygiene reading suggested (Section~\ref{sec:results:mechanism:prior_retest}) remains a reasonable move but does not change the pilot direction.
